\documentclass[11pt]{article}

\usepackage[T1]{fontenc}
\usepackage{lmodern}
\usepackage[margin=0.9in]{geometry}
\usepackage{amsmath,amssymb,mathtools}
\usepackage{booktabs,longtable,array}
\usepackage{xcolor}
\usepackage{hyperref}

\setlength{\parskip}{0.45em}
\setlength{\parindent}{0pt}
\setlength{\emergencystretch}{2em}

\hypersetup{
  colorlinks=true,
  linkcolor=blue!55!black,
  urlcolor=blue!55!black,
  pdftitle={Direct audit of the NS branching 111 coefficient},
  pdfauthor={Machine working note}
}

\newcommand{\SCA}{\mathsf{SCA}}
\newcommand{\Vir}{\mathsf{Vir}}
\newcommand{\Ff}{\mathsf F}
\newcommand{\hatrho}{\widehat\rho}
\newcommand{\rhof}{\rho^{\Ff}}
\newcommand{\half}{\tfrac12}

\title{Direct Audit of the NS Branching Coefficient in the \(111\) Channel}
\author{Machine working note}
\date{\today}

\begin{document}

\maketitle

\begin{abstract}
We audit the case \(n_1=n_2=n_3=1\) without using the blow-up product to
construct the states or to evaluate their matrix element.  First we solve for
the level-two \(\Vir\oplus\Vir\) primary \(v_1(P)\) in the auxiliary-fermion
times NS-SCA PBW basis.  Its five defining highest-weight equations and its
ungraded norm are then checked symbolically for generic \(b,P\).  We next
expand the ordered three-point coefficient into the 25 nonzero
auxiliary-fermion contractions, including the full crossing sign printed in
the human note, and evaluate every SCA descendant coefficient by Ward
identities at \((\infty,1,0)\).  Only after this direct calculation do we
compare with the four-\(\ell\) formula.  The norm agrees identically, but the
direct \(111\) three-point numerator does not agree with the printed
\(\ell\)-product.  The discrepancy is not an overall sign and survives in the
normalized square.  This note does not change the \(B^2\) used in the
double-Virasoro block.
\end{abstract}

\tableofcontents

\section{Conventions and scope}

We use exactly the parameters of the human note,
\begin{equation}
 Q=b+b^{-1},\qquad
 c=\frac32+3Q^2,\qquad
 h(P)=\frac{Q^2}{8}-\frac{P^2}{2}.
 \label{eq:parameters}
\end{equation}
The auxiliary fermion is denoted by \(f_r\) here (it is denoted by
\(\psi_r\) in the human note), and
\begin{equation}
 f_r^\dagger=-f_{-r},\qquad G_r^\dagger=G_{-r}.
 \label{eq:dagger}
\end{equation}
The tensor-product pairing is the ungraded one,
\begin{equation}
 B_{\mathrm{ug}}(u\otimes x,v\otimes y)
 =B_{\Ff}(u,v)B_{\SCA}(x,y).
 \label{eq:ungraded-pairing}
\end{equation}
The ordered holomorphic three-point form is nevertheless factorized with the
full crossing sign written in the current human note:
\begin{equation}
 \hatrho_a(u_1\otimes x_1,u_2\otimes x_2,u_3\otimes x_3)
 =(-1)^{|x_1|(|u_2|+|u_3|)+|x_2||u_3|}
 \rhof(u_1,u_2,u_3)\rho_a(x_1,x_2,x_3).
 \label{eq:full-crossing}
\end{equation}
Thus the pairing convention and the factorization convention are kept
distinct throughout the audit.

The notation ``\(111\)'' means \(n_1=n_2=n_3=1\), not the lower case
\(n_i=\half\).  The total parity is even, so the relevant form is
\(\hatrho_0\).

\section{The level-two primary \(v_1(P)\)}

\subsection{Seven-component ansatz}

At total level two and even total parity, the relevant tensor-product PBW
basis contains the following seven states:
\begin{align}
 &L_{-2}\phi,\quad L_{-1}^2\phi,\quad
 G_{-\half}G_{-\frac32}\phi,\notag\\
 &f_{-\half}G_{-\frac32}\phi,\quad
 f_{-\half}G_{-\half}L_{-1}\phi,\quad
 f_{-\frac32}G_{-\half}\phi,\quad
 f_{-\half}f_{-\frac32}\phi.
 \label{eq:seven-basis}
\end{align}
Define the abbreviations
\begin{equation}
 a=\frac Q2+P,\qquad
 d=(a+b)(a+b^{-1}),\qquad
 R=Q+P.
 \label{eq:adr}
\end{equation}
Solving
\begin{equation}
 L_m^{(j)}v_1=0\quad(j=1,2;\ m=1,2),\qquad
 (L_0^{(1)}-h_1^{(1)})v_1=0
 \label{eq:five-primary-equations}
\end{equation}
gives one direction.  Fixing its scale by the human-note Fermi-sea
normalization gives
\begin{equation}
\boxed{
\begin{aligned}
v_1(P)=\big[{}
 &4dL_{-2}-2L_{-1}^2
 -2(a^2+d)G_{-\half}G_{-\frac32}\\
 &-4R f_{-\half}G_{-\half}L_{-1}
 -4a^2R f_{-\half}G_{-\frac32}\\
 &+4dR f_{-\frac32}G_{-\half}
 -4adR f_{-\half}f_{-\frac32}
 \big]\phi(P).
\end{aligned}}
\label{eq:v1-explicit}
\end{equation}

The overall sign can be fixed without a three-point formula.  The human note
orders
\begin{equation}
 w_1=\chi_{-\frac32}\chi_{-\half}\phi,
 \qquad \chi_r=f_r-i\eta_r.
\end{equation}
In the canonical auxiliary basis \(f_{-\half}f_{-\frac32}\phi\), the
pure-\(f\) coefficient of \(w_1\) is therefore \(-1\).  Moreover,
\begin{equation}
 \ell(Q+2P,4)=16adR.
\end{equation}
Hence \(v_1=\frac14\ell(Q+2P,4)w_1\) has pure-\(f\) coefficient
\(-4adR\), exactly as in \eqref{eq:v1-explicit}.  This removes a possible
overall-sign ambiguity before the \(111\) matrix element is considered.

As an independent source check, equation (B.3a) of
\href{https://arxiv.org/pdf/1111.2803}{arXiv:1111.2803} gives the same
seven-component vector after (i) converting to auxiliary-first order,
(ii) using \(G_{-\half}^2=L_{-1}\) and
\(G_{-\frac32}G_{-\half}=2L_{-2}-G_{-\half}G_{-\frac32}\), and
(iii) multiplying their vector by \(-2\) to obtain the human-note
normalization.

\subsection{Direct primary check}

The code acts on \eqref{eq:v1-explicit} with the two commuting Virasoro
generators written in the human note.  No \(\ell\)-factor is used in these
actions.  For symbolic nonzero \(b,P\), the four positive-mode residuals and
the \(L_0^{(1)}\) residual all simplify identically to zero:
\begin{equation}
 \#\{\hbox{vanishing defining residuals}\}=5/5.
 \label{eq:primary-result}
\end{equation}

\section{Direct ungraded norm}

It is useful to split the vector according to its auxiliary state:
\begin{equation}
 v_i=Y_i^{(0)}+f_{-\half}Y_i^{(1)}
       +f_{-\frac32}Y_i^{(3)}
       +f_{-\half}f_{-\frac32}Y_i^{(13)},
 \label{eq:block-split}
\end{equation}
where, for leg \(i\), \(a,d,R\) are evaluated at \(P=P_i\), and
\begin{align}
Y^{(0)}={}&
 \big[4dL_{-2}-2L_{-1}^2
 -2(a^2+d)G_{-\half}G_{-\frac32}\big]\phi,
 \label{eq:Y0}\\
Y^{(1)}={}&
 \big[-4R G_{-\half}L_{-1}-4a^2R G_{-\frac32}\big]\phi,
 \label{eq:Y1}\\
Y^{(3)}={}&4dR G_{-\half}\phi,
\qquad
Y^{(13)}=-4adR\phi.
 \label{eq:Y313}
\end{align}
Since \(f_r^\dagger=-f_{-r}\), the four auxiliary norms have signs
\(+,-,-,+\).  The ungraded norm is therefore
\begin{equation}
 \lVert v_1\rVert_{\rm ug}^2
 =\langle Y^{(0)},Y^{(0)}\rangle
 -\langle Y^{(1)},Y^{(1)}\rangle
 -\langle Y^{(3)},Y^{(3)}\rangle
 +\langle Y^{(13)},Y^{(13)}\rangle.
 \label{eq:norm-sector-sum}
\end{equation}
The SCA Gram matrix follows only from
\begin{equation}
 [L_m,L_n]=(m-n)L_{m+n}+\frac{c}{12}(m^3-m)\delta_{m+n,0},
\end{equation}
\begin{equation}
 [L_m,G_r]=(\tfrac m2-r)G_{m+r},\qquad
 \{G_r,G_s\}=2L_{r+s}+\frac c3(r^2-\tfrac14)\delta_{r+s,0}.
\end{equation}
Exact symbolic contraction of \eqref{eq:norm-sector-sum} gives
\begin{align}
\lVert v_1(P)\rVert_{\rm ug}^2
={}&\frac{4}{b^6}P(P+b)(Pb+1)(Pb+b^2+1)
 (2Pb+b^2+1)^2\notag\\
&\hspace{1.1cm}\times
 (2Pb+b^2+3)(2Pb+3b^2+1).
 \label{eq:direct-norm-factorized}
\end{align}
On the other hand,
\begin{align}
 \ell(2P,4)&=16P(P+b)(P+b^{-1})a,\\
 \ell(Q+2P,4)&=16adR.
\end{align}
Consequently \eqref{eq:direct-norm-factorized} is exactly
\begin{equation}
 \boxed{
 \lVert v_1(P)\rVert_{\rm ug}^2
 =\frac14\ell(2P,4)\ell(Q+2P,4),}
 \label{eq:norm-agreement}
\end{equation}
which is the \(n=1\) specialization of the human-note norm formula.  This is
a generic symbolic identity, not a numerical sample.

\section{Direct \(111\) three-point coefficient}

\subsection{Fixed-parity SCA three-form}

The ordered SCA form is at \((\infty,1,0)\) and is normalized by
\begin{equation}
 \rho_0(\phi_1,\phi_2,\phi_3)=1,
 \qquad
 \rho_1(\phi_1,G_{-\half}\phi_2,\phi_3)=1.
\end{equation}
The eight global values used to terminate the Ward recursion are exactly the
two tables of the human note:
\begin{center}
\begin{tabular}{c|cccc}
\toprule
\((a_1a_2a_3)\)&000&110&101&011\\
\midrule
\(\rho_0\)&1&\(h_1+h_2-h_3\)&\(h_1-h_2+h_3\)&\(h_1-h_2-h_3\)\\
\bottomrule
\end{tabular}

\vspace{0.5em}
\begin{tabular}{c|cccc}
\toprule
\((a_1a_2a_3)\)&100&010&001&111\\
\midrule
\(\rho_1\)&1&1&\(-1\)&\(-(h_1+h_2+h_3-\half)\)\\
\bottomrule
\end{tabular}
\end{center}
Here \(a_i\) is the power of \(G_{-\half}\) on leg \(i\).  In the code this
fixed-parity table is obtained from the internal Ward-engine component by the
conversion
\begin{equation}
 \rho_a^{\rm fixed}(x_1,x_2,x_3)
 =(-1)^{a|x_3|}\rho^{\rm engine}(x_1,x_2,x_3),
 \qquad a=|x_1|+|x_2|+|x_3|\pmod2.
 \label{eq:fixed-conversion}
\end{equation}
This conversion is tested against all eight entries above before any level-two
coefficient is evaluated.

All remaining SCA descendants in \eqref{eq:Y0}--\eqref{eq:Y313} are reduced
recursively by the NS Ward identities.  For example, if \(B,M,K\) denote the
three states, then
\begin{equation}
 \rho(B,L_{-1}M,K)
 =(\Delta_B-\Delta_M-\Delta_K)\rho(B,M,K),
 \label{eq:ward-L1}
\end{equation}
and, for \(n\ge2\),
\begin{align}
\rho(B,L_{-n}M,K)
={}&\sum_{j\ge0}\binom{n-2+j}{n-2}
 \rho(L_{n+j}B,M,K)\notag\\
&+(-1)^n\sum_{j\ge0}\binom{n-2+j}{n-2}
 \rho(B,M,L_{j-1}K).
 \label{eq:ward-Ln}
\end{align}
For an NS mode \(r\in\mathbb Z+\half\), \(r>0\), the middle-leg recursion is
\begin{align}
\rho(B,G_{-r}M,K)
={}&\sum_{j\ge0}\binom{r-\frac32+j}{j}
 \rho(G_{r+j}B,M,K)\notag\\
&-(-1)^{|B|+|K|+r+\half}
 \sum_{j\ge0}\binom{r-\frac32+j}{j}
 \rho(B,M,G_{j-\half}K).
 \label{eq:ward-Gr}
\end{align}
The sums terminate because the external modules are highest-weight modules.
Bra modes are moved to the middle and ket legs with the corresponding contour
identity; ket modes are reduced by endpoint reflection.  Thus every term in
the calculation is obtained from the displayed algebra and the eight base
values, with no input from the \(\ell\)-product.

\subsection{Auxiliary-fermion contractions and crossing signs}

Use the labels
\begin{equation}
 0=\mathbf1_{\Ff},\qquad
 1=f_{-\half},\qquad
 3=f_{-\frac32},\qquad
 13=f_{-\half}f_{-\frac32}.
\end{equation}
Every summand of \(v_1\) is even in total, so its SCA parity equals its
auxiliary parity.  Therefore the exponent in \eqref{eq:full-crossing} reduces
for the \(111\) expansion to
\begin{equation}
 E=p_1p_2+p_1p_3+p_2p_3,
 \qquad p_i=|u_i|.
 \label{eq:reduced-crossing}
\end{equation}
The auxiliary three-form is evaluated directly as the Pfaffian of the sphere
fermion two-point function.  The bra word is reversed, and the rule
\(f_r^\dagger=-f_{-r}\) supplies one minus sign per bra fermion.  Exactly 25
triples survive:

\begingroup\small
\begin{longtable}{ccc|rrr@{\qquad}ccc|rrr}
\toprule
\(A\)&\(B\)&\(C\)&\(\rhof\)&\(E\)&\((-1)^E\rhof\)&
\(A\)&\(B\)&\(C\)&\(\rhof\)&\(E\)&\((-1)^E\rhof\)\\
\midrule
\endhead
0&0&0&1&0&1 & 0&1&1&1&1&-1\\
0&1&3&1&1&-1 & 0&3&1&-1&1&1\\
0&3&3&-2&1&2 & 0&13&13&1&0&1\\
1&0&1&-1&1&1 & 1&1&0&-1&1&1\\
1&1&13&1&1&-1 & 1&3&13&-2&1&2\\
1&13&1&1&1&-1 & 1&13&3&2&1&-2\\
3&0&3&-1&1&1 & 3&1&0&-1&1&1\\
3&1&13&-1&1&1 & 3&3&0&-1&1&1\\
3&3&13&1&1&-1 & 3&13&1&2&1&-2\\
3&13&3&3&1&-3 & 13&0&13&1&0&1\\
13&1&1&-1&1&1 & 13&1&3&1&1&-1\\
13&3&1&-1&1&1 & 13&13&0&1&0&1\\
13&13&13&-4&0&-4 & & & & & \\
\bottomrule
\end{longtable}
\endgroup

Writing \(Y_{i}^{(A)}\) for the SCA state multiplying auxiliary state \(A\)
on leg \(i\), the direct result is therefore the finite sum
\begin{equation}
 \boxed{
 T_{111}^{\rm dir}
 =\sum_{(A,B,C)\ \text{in the table}}
 (-1)^{E_{ABC}}\rhof(A,B,C)
 \rho_0(Y_1^{(A)},Y_2^{(B)},Y_3^{(C)}).}
 \label{eq:direct-25-sum}
\end{equation}
The implementation also evaluates the seven-by-seven-by-seven expansion
without using this table.  The unsplit expansion and the 25-block sum agree
exactly; this is an internal check on the bookkeeping in
\eqref{eq:direct-25-sum}.

\section{Comparison with the human-note \(\ell\)-product}

For \(n_1=n_2=n_3=1\), the printed numerator specializes to
\begin{align}
T_{111}^{\ell}
={}&-\frac18
 \ell\left(\frac Q2+P_1+P_2+P_3,6\right)
 \ell\left(\frac Q2-P_1+P_2+P_3,2\right)\notag\\
&\quad\times
 \ell\left(\frac Q2+P_1+P_2-P_3,2\right)
 \ell\left(\frac Q2-P_1+P_2-P_3,-2\right).
 \label{eq:ell-111-raw}
\end{align}
Since
\begin{equation}
 \ell(x,2)=x,\qquad \ell(x,-2)=-\ell(Q-x,2)=-(Q-x),
\end{equation}
this becomes
\begin{equation}
\boxed{
T_{111}^{\ell}=\frac18\ell(S,6)ABC,}
\label{eq:ell-111-reduced}
\end{equation}
where
\begin{align}
S&=\frac Q2+P_1+P_2+P_3,&
A&=\frac Q2-P_1+P_2+P_3,\\
B&=\frac Q2+P_1+P_2-P_3,&
C&=\frac Q2+P_1-P_2+P_3.
\end{align}

\subsection{Exact one-parameter comparison}

To keep the symbolic output readable, take the nonsingular line
\begin{equation}
 b=3,\qquad (P_1,P_2,P_3)=(-\tfrac32,t,-\tfrac32).
 \label{eq:compact-line}
\end{equation}
The direct Ward calculation gives
\begin{equation}
T_{111}^{\rm dir}
=-\frac{(3t+5)}{4251528}\,
\begin{aligned}[t]
\big(&177147t^{11}+885735t^{10}-27359370t^9-347601780t^8\\
&-1767367188t^7-4991980590t^6-7102586493t^5\\
&+1135874745t^4+12286758429t^3+1003438035t^2\\
&+6876067671t+31288240715\big).
\end{aligned}
\label{eq:direct-line}
\end{equation}
The human-note product instead gives the completely split polynomial
\begin{align}
T_{111}^{\ell}
={}&-\frac{1}{157464}
t(t+2)(t+8)(3t-4)(3t-2)(3t+4)\notag\\
&\qquad\times(3t+5)^2(3t+8)(3t+14)(3t+16)(3t+32).
\label{eq:ell-line}
\end{align}
Equations \eqref{eq:direct-line} and \eqref{eq:ell-line} are different
polynomials.  In particular, their difference is not removed by an overall
sign or phase.

At the exact rational point \(t=\frac13\), all three norms are nonzero and
the comparison is
\begin{align}
T_{111}^{\rm dir}&=-\frac{315962090}{6561},&
T_{111}^{\ell}&=-\frac{818125}{486},\\
T_{111}^{\rm dir}-T_{111}^{\ell}
&=-\frac{609834805}{13122},&
\frac{T_{111}^{\rm dir}}{T_{111}^{\ell}}
&=\frac{126384836}{4417875}.
\label{eq:sample-rho}
\end{align}
The three independently verified norms are
\begin{equation}
 N_1=N_3=\frac{1463}{108},\qquad
 N_2=\frac{1971200}{243}.
\end{equation}
Thus the normalized squares also differ:
\begin{align}
\frac{(T_{111}^{\rm dir})^2}{N_1N_2N_3}
&=\frac{998320423171681}{640775109744},\\
\frac{(T_{111}^{\ell})^2}{N_1N_2N_3}
&=\frac{13546875}{7116032}.
 \label{eq:sample-B2}
\end{align}
The second value in \eqref{eq:sample-B2} is the existing blow-up \(B^2\)
used by the double-Virasoro block.  It has not been modified by this audit.

As a second generic exact sample, with
\begin{equation}
 b=\frac54,\qquad(P_1,P_2,P_3)=\left(\frac16,\frac27,-\frac38\right),
\end{equation}
the ratio is
\begin{equation}
 \frac{T_{111}^{\rm dir}}{T_{111}^{\ell}}
 =\frac{40988240919614376360207863310856861}
 {3383638222509524104324824472946800}\ne1.
\end{equation}
All arithmetic in both samples is exact rational arithmetic.

\section{Conclusion and interpretation}

The extent of the check is:
\begin{enumerate}
\item The explicit seven-component \(v_1(P)\) satisfies all five defining
  primary/eigenvalue equations symbolically for generic \(b,P\).
\item Its ungraded norm agrees symbolically with the \(n=1\) specialization
  of the human-note norm formula.
\item All 25 nonzero auxiliary-fermion triples are included with
  \(f_r^\dagger=-f_{-r}\) and the full crossing exponent in
  \eqref{eq:full-crossing}.
\item Every SCA coefficient is reduced to the fixed-parity base table by
  exact Ward identities.  The split 25-term sum agrees with a separate
  unsplit expansion.
\item The direct result disagrees with the printed four-\(\ell\) numerator on
  an exact one-parameter family and at two explicit rational samples.  The
  disagreement survives in \(B^2\).
\end{enumerate}

Therefore the following set of statements is not simultaneously consistent:
the ungraded tensor pairing, \(f_r^\dagger=-f_{-r}\), the full signed
factorization \eqref{eq:full-crossing}, the fixed-parity SCA base table, the
normalization \eqref{eq:v1-explicit}, and the identification of the resulting
ordered three-point form with the printed \(\ell\)-product.  The primary and
norm checks isolate the remaining issue to the dictionary between the
ordered \(\hatrho\) in \eqref{eq:full-crossing} and the convention in which
the blow-up product was derived.  A single additional overall sign cannot
repair the \(111\) result.

This conclusion is deliberately narrower than changing the production
branching coefficient.  The existing \(B^2\) used in the double-Virasoro
block is left intact because it is separately constrained by the block-level
comparison with NS \(c\)-recursion.  What must be repaired is the convention
transport that claims that this \(B^2\) is obtained by directly squaring the
specific ordered \(\hatrho\) defined by the simultaneous conventions above.

\section{Reproduction}

The exact calculation is implemented in
\begin{center}
\texttt{Code/audit\_ns\_branching\_111.py}.
\end{center}
Run from the repository root:
\begin{verbatim}
python3 Code/audit_ns_branching_111.py --compact-line
\end{verbatim}
The script constructs \eqref{eq:v1-explicit} directly, performs the generic
primary and norm checks, prints the complete auxiliary table, evaluates the
two rational samples, and factors the one-parameter polynomials
\eqref{eq:direct-line} and \eqref{eq:ell-line}.  It imports the existing Ward
and Gram engines but does not write to, replace, or patch the production
double-Virasoro branching coefficient.

\end{document}
